\problem{6.32}{SU(2) rotations of a spin-half particle}
Let $\boldsymbol\sigma=(\sigma_x,\sigma_y,\sigma_z)$
be the Pauli matrices.
\textbf{(a)} Their multiplication rule is
\[
\sigma_i\sigma_j=\delta_{ij}I+
\ii\epsilon_{ijk}\sigma_k.
\]
For ordinary vectors $\mathbf a,\mathbf b$,
sum $a_ib_j\sigma_i\sigma_j$ to obtain
\begin{equation}
\boxed{(\mathbf a\cdot\boldsymbol\sigma)
(\mathbf b\cdot\boldsymbol\sigma)
=(\mathbf a\cdot\mathbf b)I+
\ii(\mathbf a\times\mathbf b)\cdot\boldsymbol\sigma.}
\end{equation}
\textbf{(b)} For a unit vector $\mathbf n$,
the previous formula gives
$(\mathbf n\cdot\boldsymbol\sigma)^2=I$.
Thus its even powers are $I$ and odd
powers are $\mathbf n\cdot\boldsymbol\sigma$.
Expand the exponential:
\begin{align}
R_{\mathbf n}(\varphi)
&=\exp[-\ii\varphi
\mathbf n\cdot\mathbf S/\hb]\\
&=\exp[-\ii(\varphi/2)
\mathbf n\cdot\boldsymbol\sigma]\\
&=\boxed{\cos(\varphi/2)I-
\ii\sin(\varphi/2)
\,\mathbf n\cdot\boldsymbol\sigma.}
\end{align}
\textbf{(c)} With
$\mathbf n=(\sin\theta\cos\phi,
\sin\theta\sin\phi,\cos\theta)$,
\[
\mathbf n\cdot\boldsymbol\sigma=
\begin{pmatrix}
\cos\theta&\sin\theta\,\ee^{-\ii\phi}\\
\sin\theta\,\ee^{\ii\phi}&-\cos\theta
\end{pmatrix}.
\]
Substitute this matrix into part (b)
to obtain the requested standard-basis
$2\times2$ rotation matrix.
\textbf{(d)} Let $N=\mathbf n\cdot\boldsymbol\sigma$,
which is Hermitian and satisfies $N^2=I$.
Then
\[
R^\dagger R=(cI+\ii sN)(cI-\ii sN)
=(c^2+s^2)I=I
\]
with $c=\cos(\varphi/2)$,
$s=\sin(\varphi/2)$.
Thus $\boxed{R^\dagger R=I}$.
\textbf{(e)} For a $z$ rotation,
$R_z=\operatorname{diag}(\ee^{-\ii\varphi/2},
\ee^{+\ii\varphi/2})$. Multiply explicitly:
\begin{align}
R_z^\dagger S_xR_z
&=\frac{\hb}{2}
\begin{pmatrix}0&\ee^{\ii\varphi}\\
\ee^{-\ii\varphi}&0\end{pmatrix}\\
&=\boxed{S_x\cos\varphi-S_y\sin\varphi}.
\end{align}
This agrees with the transformed
vector component in the textbook's
active-state convention.
\textbf{(f)} Around $x$ through $\pi$,
\[
\boxed{R_x(\pi)=-\ii\sigma_x=
\begin{pmatrix}0&-\ii\\-\ii&0\end{pmatrix}.}
\]
It sends $\ket{\uparrow_z}=(1,0)^T$
to $-\ii\ket{\downarrow_z}$, so the
spin is flipped up to an
unobservable overall phase.
\textbf{(g)} A complete $2\pi$ turn gives
\begin{equation}
\boxed{R_z(2\pi)=
\operatorname{diag}(-1,-1)=-I.}
\end{equation}
A spinor changes sign under $2\pi$
rotation; it returns to itself only
after $4\pi$. The sign is physically
observable through interference
against an unrotated reference path.

\problem{6.33}{Dihedral symmetry and a degenerate square-well pair}
Use the centered square $-L/2<x,y<L/2$,
and write
$f_n(x)=\sqrt{2/L}\sin[n\pi(x/L-1/2)]$,
$\psi_{ab}(x,y)=f_a(x)f_b(y)$.
\textbf{(a)} Under active counterclockwise
rotation by $\pi/2$, the wavefunction
is evaluated at the inverse-rotated
coordinates $(y,-x)$:
\[
(\RR_{\pi/2}\psi_{ab})(x,y)
=f_a(y)f_b(-x).
\]
An individual centered box mode
satisfies
\[
f_n(-x)=(-1)^{n+1}f_n(x).
\]
Therefore
\begin{equation}
\boxed{\RR_{\pi/2}\psi_{ab}
=(-1)^{b+1}\psi_{ba},\qquad
\RR_{\pi/2}\psi_{ba}
=(-1)^{a+1}\psi_{ab}.}
\end{equation}
\textbf{(b)} If $a$ and $b$ have
the same parity, then
$s=(-1)^{a+1}=(-1)^{b+1}$.
Define $\psi_\pm=(\psi_{ab}\pm\psi_{ba})/\sqrt2$.
Using part (a),
\[
\RR_{\pi/2}\psi_\pm=
s(\psi_{ba}\pm\psi_{ab})/\sqrt2,
\]
so
\begin{equation}
\boxed{\RR_{\pi/2}\psi_+=s\psi_+,\qquad
\RR_{\pi/2}\psi_-=-s\psi_-.}
\end{equation}
For two odd indices, $s=+1$;
for two even indices, $s=-1$.
\textbf{(c)} For $a=5$, $b=7$,
both one-dimensional modes are even
under $x\mapsto-x$.
Hence reflections in the
horizontal or vertical axis preserve
$\psi_-$. Reflection in the
diagonal $x=y$ interchanges the
two products and changes the sign:
\[
\psi_-(y,x)=-\psi_-(x,y).
\]
The other diagonal combines a
sign-preserving axial reflection
with this interchange, so it also
gives a minus sign.
A quarter turn gives a minus sign
by (b), and its powers produce
$\pm1$ factors. Therefore $\psi_-$
is an eigenstate of \emph{every}
symmetry in the square's dihedral
group $D_4$.
The script writes a contour plot of
$\psi_-=\frac1{\sqrt2}(\psi_{57}-\psi_{75})$,
whose nodal and sign pattern visibly
obeys these transformations.
\phys{Although $\psi_{57}$ and
$\psi_{75}$ have the same energy
$E\propto5^2+7^2$, their symmetric
and antisymmetric combinations
transform in inequivalent one-dimensional
$D_4$ representations. This discrete
symmetry alone does not \emph{force}
their equality of energy; the
separable square Hamiltonian has
additional structure.}

\problem{6.34}{Laplace--Runge--Lenz operator and hidden hydrogen degeneracy}
Put $k=e^2/(4\pi\epsilon_0)$.
The Hermitian quantum
Laplace--Runge--Lenz vector is
\[
\mathbf M=\frac{\mathbf p\times\mathbf L-
\mathbf L\times\mathbf p}{2m}
-\frac{k}{r}\,\mathbf r.
\]
The exercise supplies the algebra
$[H,M_i]=0$, $[L_i,M_j]=
\ii\hb\epsilon_{ijk}M_k$,
$[M_i,M_j]=-(2\ii\hb H/m)
\epsilon_{ijk}L_k$, together with
\begin{equation}
\mathbf M^2=k^2+\frac{2H}{m}
(\mathbf L^2+\hb^2),
\qquad\mathbf M\cdot\mathbf L=0.
\end{equation}
The latter are taken as given,
as specified by the exercise.
Write $\ket{n\ell\ell}$ for a
hydrogen state with maximal $m$.
Rotational covariance (Problem 6.19)
and $[H,M_+]=0$ imply
\[
M_+\ket{n\ell\ell}
=c_{n\ell}\ket{n,\ell+1,\ell+1}
\]
unless the state is annihilated.
\textbf{(a)} The identity
$\mathbf M\cdot\mathbf L=0$ can be
written
\[
M_zL_z+\frac12(M_+L_-+M_-L_+)=0.
\]
Apply it to the highest-weight
state, for which $L_+\ket{n\ell\ell}=0$.
Use $L_-\ket{\ell\ell}
=\hb\sqrt{2\ell}\ket{\ell,\ell-1}$.
The Wigner--Eckart ratio of relevant
rank-one tensor matrix elements gives
\[
M_+\ket{n\ell,\ell-1}
=c_{n\ell}\sqrt{\frac{\ell}{\ell+1}}
\ket{n,\ell+1,\ell},
\]
which follows from the ordinary
CG coefficients for $\ell\otimes1
\to\ell+1$.
For $\ell>0$ substitute to obtain
\begin{equation}
\boxed{M_z\ket{n\ell\ell}
=-\frac{c_{n\ell}}{\sqrt{2(\ell+1)}}
\ket{n,\ell+1,\ell}.}
\end{equation}
The $\ell=0$ case follows by
rotational covariance / CG factors,
since the scalar product relation
alone becomes trivial there.
\textbf{(b)} With
$M_\pm=M_x\pm\ii M_y$,
\[
M_-M_+=M_x^2+M_y^2+
\ii[M_x,M_y]
=M^2-M_z^2+\frac{2\hb H}{m}L_z.
\]
Act on $\ket{n\ell\ell}$ and insert
$E_n=-mk^2/(2\hb^2n^2)$,
$\mathbf L^2=\hb^2\ell(\ell+1)$ and
$L_z=\hb\ell$:
\begin{align}
M_-M_+\ket{n\ell\ell}
&=\left[k^2-\frac{k^2}{n^2}
(\ell(\ell+1)+1)
-\frac{k^2\ell}{n^2}-M_z^2\right]
\ket{n\ell\ell}\\
&=\boxed{\left[
k^2\left(1-\frac{(\ell+1)^2}{n^2}\right)
-M_z^2\right]\ket{n\ell\ell}.}
\end{align}
\textbf{(c)} Take the expectation
value in the normalized state.
Because $M_-=M_+^\dagger$,
the left side is
$\|M_+\ket{n\ell\ell}\|^2=|c_{n\ell}|^2$.
From part (a),
$\|M_z\ket{n\ell\ell}\|^2=
|c_{n\ell}|^2/[2(\ell+1)]$.
Therefore
\[
|c_{n\ell}|^2=
k^2\left(1-\frac{(\ell+1)^2}{n^2}\right)
-\frac{|c_{n\ell}|^2}{2(\ell+1)}.
\]
Solving for the coefficient:
\begin{equation}
\boxed{|c_{n\ell}|=
k\sqrt{\frac{2(\ell+1)}{2\ell+3}
\left[1-\frac{(\ell+1)^2}{n^2}\right]}.}
\end{equation}
An arbitrary phase may be chosen
for the states, so only the magnitude
is physically fixed.
This coefficient is nonzero for
$0\le\ell\le n-2$, but vanishes
for $\ell=n-1$, where no state
with $\ell=n$ exists in the same
hydrogen shell.
\phys{The LRL vector moves between
different $\ell$ multiplets at fixed
principal quantum number $n$.
Together with angular momentum,
it generates the hidden SO(4)
symmetry behind the $n^2$
orbital degeneracy of hydrogen
(with spin ignored).}
